\documentclass[10pt,a4paper]{amsart}
\usepackage[margin=19mm]{geometry}
\usepackage{amsmath,amssymb,mathtools}
\usepackage[scr=rsfs,cal=euler]{mathalfa}
\usepackage{xfrac}
\usepackage[unicode,hidelinks,bookmarks=false]{hyperref}
\newcommand{\ie}{i.e.\ }
\newcommand{\cf}{cf.\ }
\newcommand{\resp}{respectively\ }
\newcommand{\Subsection}{\subsection}
\providecommand{\nobreakdash}{\nobreak\hbox{-}}
\setlength{\emergencystretch}{2em}
\multlinegap0pt
\usepackage{amssymb}
\usepackage{tikz}
\usepackage{dsfont}
\usepackage[shortlabels]{enumitem}
\usepackage{float}
\newtheorem{lemma}{Lemma}
\newtheorem{empheq}{Empheq}
\title{Backward problem for a degenerate viscous Hamilton-Jacobi equation: stability and numerical identification}
\author{S. E. Chorfi, A. Habbal, M. Jahid, L. Maniar, A. Ratnani}
\date{Source: arXiv:2603.09509v1 (CC BY 4.0)}
\begin{document}
\maketitle

\noindent\emph{Source and licence.} Backward problem for a degenerate viscous Hamilton-Jacobi equation: stability and numerical identification. Version arXiv:2603.09509v1 (CC BY 4.0), distributed by arXiv under the Creative Commons Attribution 4.0 International licence, http://creativecommons.org/licenses/by/4.0/, as recorded in the arXiv licence metadata for this exact version and shown on the abstract page https://arxiv.org/abs/2603.09509v1. Full licence notice: \texttt{licenses/arxiv-2603.09509-CC-BY-4.0.txt}.\par\smallskip

\noindent\textit{Editorial scope.} Counted unit: the article's lemma lemma (source0.tex lines 751--754). The article's complete original proof of it is delivered with it (source0.tex lines 755--803, one \texttt{proof} environment of 49 lines). No mathematics is written, repaired or re-derived: the statement and the proof are the article's own lines, in the article's own notation, and no retained line is substituted or re-flowed.  No further source-local context is needed: every cross-reference of every retained line resolves inside the two delivered spans.  Nothing was omitted from any retained span, so the package carries no omission record.  No retained line contains a citation, so the wrapper carries no bibliography and no bibliography entry.  No retained line contains a figure environment, an image-inclusion command or drawing code, so no image is delivered and the package declares no asset dependency.  Editorial formatting only, in addition: an AMS \texttt{amsart} preamble that restates the article's own declarations and macros used by the retained lines, namely the environments \texttt{lemma}, \texttt{empheq} on separate counters, together with the standard packages those lines need.  The retained environments print in this document as Lemma 1, Empheq 1 in order of appearance, whereas the article prints the same items with the numbering of its own full document.  The complete native source of the article as distributed in the arXiv e-print for this exact version is bundled unchanged as \texttt{sources/196067/source0.tex}.
\begin{lemma}
The Fréchet derivative $J'$ is Lipschitz continuous with Lipschitz constant $L=1$, i.e.,
    $$ \vert \vert J'(f + \delta f )-J'(f)\vert\vert_{\frac{1}{a}}\leqslant  \vert \vert \delta f\vert \vert_{\frac{1}{a}} \qquad \forall f,\, f+\delta f \in W.$$
\end{lemma}

\begin{proof}
    By definition 
    $$\vert\vert J'(f+\delta f)-J'(f)\vert\vert_{\frac{1}{a}}=\vert\vert\delta \psi(\cdot,0;f)\vert\vert_{\frac{1}{a}},$$
    where
    $\delta\psi(x,t;f)= \psi(x,t;f+\delta f)-\psi(x,t;f)$ solves
%\begin{empheq}[left = \empheqlbrace]{alignat=2}
\begin{equation}
    \left\{
\begin{aligned}
& \delta\psi_{t}(x,t) = -a(x) \delta\psi_{xx}(x,t),&&\quad \text { in } Q,\\
& \delta\psi(x,t)=0, &&\quad \text { on } \Sigma,\\
& \delta\psi(x,T)=\delta u(x,T;f), &&\quad \text { in } \Omega,
\label{eq1to7}
\end{aligned}
\right.
\end{equation}
%\end{empheq}
and $\delta u(x,t;f)$ is the weak solution of the auxiliary problem
%\begin{empheq}[left = \empheqlbrace]{alignat=2}
\begin{equation}
    \left\{
\begin{aligned}
& \delta u_{t}(x,t) = a(x) \delta u_{xx}(x,t),&&\quad \text { in } Q,\\
& \delta u(x,t)=0, &&\quad \text { on } \Sigma,\\
& \delta u(x,0)=\delta f(x), &&\quad \text { in } \Omega.
\label{eq1to8}
\end{aligned}
\right.
\end{equation}
%\end{empheq}
We use the energy identity
\begin{equation*}
    \begin{aligned}
        \frac{1}{2} \frac{d}{dt}\int_{0}^1 \frac{1}{a(x)}\left(\delta \psi(x,t)\right)^2 dx&= \int_{0}^1 \frac{1}{a(x)}\delta \psi \delta \psi_{t} dx\\& =\int_{0}^1 \frac{1}{a(x)}\delta \psi a(x) \delta \psi_{xx} dx\\&= -\int_{0}^1 \delta \psi \delta \psi_{xx}\\&= \int_{0}^1 (\delta \psi_{x})^2 dx \geqslant 0.
    \end{aligned}
\end{equation*}
Integrating on $[0,T]$ and using $\delta \psi(x,T;f)=\delta u(x,T;f)$, we obtain
$$\int_{0}^1 \frac{1}{a(x)}\left(\delta \psi(x,0;f)\right)^2 dx \leqslant \int_{0}^1 \frac{1}{a(x)}\left(\delta \psi(x,T;f)\right)^2 dx=  \int_{0}^1 \frac{1}{a(x)}\left(\delta u(x,T;f)\right)^2 dx.$$
Next, we employ the following equality
\begin{equation*}
  \begin{aligned}
        \frac{1}{2} \frac{d}{dt}\int_{0}^1 \frac{1}{a(x)}\left(\delta u(x,t)\right)^2 dx&= \int_{0}^1 \frac{1}{a(x)}\delta u \delta u_{t} dx\\& =\int_{0}^1 \frac{1}{a(x)}\delta u a(x) \delta u_{xx} dx\\&= \int_{0}^1 \delta u\delta u_{xx}\\&=- \int_{0}^1 (\delta u_{x})^2 dx \leqslant 0.
  \end{aligned}  
\end{equation*}
Integrating on $[0,T]$, we obtain 
$$\int_{0}^1 \frac{1}{a(x)}\left(\delta u(x,T;f)\right)^2 dx \leqslant \int_{0}^1 \frac{1}{a(x)}\left(\delta u(x,0;f)\right)^2 dx= \int_{0}^1 \frac{1}{a(x)}\vert\delta f\vert^2 dx.$$
The above estimates yield
$$\vert\vert\delta \psi(\cdot,0;f)\vert\vert_{\frac{1}{a}}\leqslant  \vert\vert\delta f\vert\vert_{\frac{1}{a}}.$$
\end{proof}

\end{document}
